\chapter{What an observer can recover}\label{ch:access}

The previous chapters studied the topology of the history complex $K_A$ directly. An observer of a generative system rarely sees $K_A$; she sees a count of histories, the degrees of the swap graph, the sizes of neighbourhoods. An \emph{observation} is a function of the admissible set $A$ (equivalently of the regime $(\lambda,m)$), and recovery questions are questions about which functions of $A$ are functions of which others. This chapter fixes the observation ladder, proves the exact relations among its levels, develops the accessibility profile of the swap graph, and proves the radius-two identity that governs everything an observer learns from two-step neighbourhoods. Whether these observations determine \emph{topology} is the subject of \cref{ch:separation}.

\section{The observation ladder}\label{sec:ladder}
Throughout, $k=6$ unless stated, $A\subseteq S_k$ is an admissible set, $G_A$ its legal-swap graph (\cref{def:swap}) and $K_A$ the permutohedral complex (\cref{def:KA}). Write $d(w)$ for the degree of $w$ in $G_A$, $B_r(w)$ for the graph ball of radius $r$ and $N=|A|$.
@@def:ladder@@

\begin{figure}[ht]\centering
\begin{tikzpicture}[font=\small,
 okn/.style={minimum width=1.9cm,minimum height=0.75cm,draw,fill=accent!16,align=center},
 badn/.style={minimum width=1.9cm,minimum height=0.75cm,draw,fill=red!12,align=center},
 unkn/.style={minimum width=1.9cm,minimum height=0.75cm,draw,fill=gray!10,align=center}]
 \foreach \i/\l in {0/{$O_N$},1/{$O_h$},2/{$O_G$},3/{$O_C$},4/{$O_A$}}{\node[font=\bfseries] at (2.0*\i,1.5) {\l};}
 \node[anchor=east] at (-1.1,0.8) {$\beta_0$};
 \node[anchor=east] at (-1.1,0.0) {$\beta_1$};
 \node[anchor=east] at (-1.1,-0.8) {$\beta_2>0$};
 \node[anchor=east] at (-1.1,-1.6) {anchored state};
 \node[unkn] at (0,0.8) {not claimed};\node[unkn] at (2,0.8) {not claimed};
 \foreach \i in {2,3,4}{\node[okn] at (2.0*\i,0.8) {determines};}
 \node[unkn] at (0,0.0) {not claimed};\node[unkn] at (2,0.0) {not claimed};
 \foreach \i in {2,3,4}{\node[okn] at (2.0*\i,0.0) {determines$^\ast$};}
 \foreach \i in {0,1}{\node[badn] at (2.0*\i,-0.8) {fails};}
\node[unkn] at (4.0,-0.8) {open};
 \foreach \i in {3,4}{\node[okn] at (2.0*\i,-0.8) {determines};}
 \foreach \i in {0,1,2}{\node[badn] at (2.0*\i,-1.6) {fails};}
\node[unkn] at (6.0,-1.6) {open};
 \node[okn] at (8,-1.6) {determines};
\end{tikzpicture}
\caption{The observation ladder against four targets. Blue: the level determines the target on all admissible sets; red: distinct admissible sets with equal observation and different target exist; grey (open): neither is claimed. $^\ast$For prefix-defined $A$ and $k\le8$ (\cref{thm:skeleton}). For $\beta_2>0$ the uncoloured graph fixes everything but $\operatorname{rank}\partial_3$ (\cref{thm:third-boundary}); whether it determines $\beta_2$ on realizable systems is not settled here. For the anchored state, failure at $O_G$ (\cref{prop:uncoloured-fails}) implies failure at the coarser levels. Whether radius-two data $O_h$ still recover $\beta_2>0$ on the arithmetic census is a different question, answered in \cref{ch:separation}.}\label{fig:ladder}
\end{figure}

@@thm:obs-ladder+analytic@@
\subsection*{The anchored state is not a function of uncoloured data}
For the two anchor families (\cref{ch:topology}) define the \emph{anchored hole state} $Y_{\mathrm{anc}}(m)\in\{1,0,\mathrm{NA}\}$: $1$ if the generator of $H_2(K_A(m_0))$ maps to a nonzero class in $H_2(K_A(m))$, $0$ if it maps to zero, and $\mathrm{NA}$ if $A(m)\not\supseteq A(m_0)$ so that the map is undefined.
@@prop:uncoloured-fails+finite-exact@@
The point is conceptual: the anchored state depends on the identification of a specific cycle across sources, which is labelled information, whereas $\mathbf 1[\beta_2>0]$ depends only on the complex.

\section{The accessibility profile}\label{sec:profile}
Define the \emph{recombination information at radius $r$} for the uniform measure on $A$:
\[
I_r=\frac1N\sum_{w\in A}\log_2|B_r(w)|=\frac1N\log_2P_r,\qquad P_r=\prod_{w\in A}|B_r(w)|\in\mathbb N .
\]
$I_r$ is the mean of the Hartley entropy of the number of histories reachable from $w$ in at most $r$ swaps.
@@thm:profile+analytic@@
\begin{remark}
Statement (e) says that sufficiently large radius-one information \emph{certifies some graph cycle}. It does not certify a hole in the complex: graph cycles are filled by squares and hexagons, and whether a cycle is filled is the content of \cref{ch:topology,ch:atlas}.
\end{remark}

\section{The radius-two identity}\label{sec:radius2}
Define for $w\in A$ the quantities $s(w)=\sum_{u\sim w}(d(u)-1)$ (the number of non-backtracking two-step walks starting at $w$) and $q(w)$, the number of square faces of $K_A$ through $w$. Let $E$ be the number of edges and $Q=\tfrac14\sum_wq(w)$ the number of square faces.
@@thm:radius2+analytic@@
@@cor:moments+analytic@@
\begin{remark}[what the identity does and does not say]
The identity shows that fine radius-two data determine $N$, $E$ and the number $Q$ of commuting squares, and contain no \emph{direct} term for hexagonal (braid) faces, whose boundary $6$-cycles first affect balls of radius three. This is an interpretation: $s(w)$ and $q(w)$ are graph quantities whose joint distribution could correlate with hexagons. The witnesses of \cref{prop:witnesses} show that any such correlation does not recover topology in general, while on the arithmetic census radius-two data do separate the hole indicator (\cref{obs:residual}); why is open.
\end{remark}
